\documentclass[11pt,a4paper]{article}

\usepackage[top=2cm,bottom=2cm,left=2cm,right=2cm]{geometry}
\usepackage[spanish]{babel}
\usepackage[utf8]{inputenc}
\usepackage[T1]{fontenc}
\usepackage{amsmath,amssymb,amsthm}
\usepackage{mathtools}
\usepackage{booktabs}
\usepackage{enumitem}
\usepackage{xcolor}
\usepackage{array}
\usepackage{multirow}
\usepackage{microtype}
\usepackage[most]{tcolorbox}
\usepackage{tikz}
\usetikzlibrary{arrows.meta,positioning,shapes.geometric}
\usepackage{hyperref}

%-- Colores -------------------------------------------------------------------
\definecolor{mainblue}{RGB}{0,70,127}
\definecolor{darkgray}{RGB}{60,60,60}
\definecolor{lightgray}{RGB}{240,240,240}
\definecolor{accentgreen}{RGB}{0,120,80}
\definecolor{warnorange}{RGB}{180,80,0}

%-- Cajas personalizadas ------------------------------------------------------
\tcbuselibrary{skins,breakable,theorems}

\newtcolorbox{definicion}[1]{
  enhanced, breakable,
  colback=mainblue!6, colframe=mainblue,
  fonttitle=\bfseries\small, title={#1},
  left=4pt, right=4pt, top=3pt, bottom=3pt
}

\newtcolorbox{ojo}{
  enhanced, breakable,
  colback=warnorange!8, colframe=warnorange,
  fonttitle=\bfseries\small, title={!`OJO!},
  left=4pt, right=4pt, top=3pt, bottom=3pt
}

\newtcolorbox{destacado}{
  enhanced, breakable,
  colback=accentgreen!7, colframe=accentgreen,
  left=4pt, right=4pt, top=3pt, bottom=3pt
}

\newtcolorbox{examen}{
  enhanced, breakable,
  colback=lightgray, colframe=darkgray,
  fonttitle=\bfseries\small, title={[Examen]},
  left=4pt, right=4pt, top=3pt, bottom=3pt
}

%-- Macro diagrama ------------------------------------------------------------
\newcommand{\diagcaption}[1]{\begin{center}\small\textit{#1}\end{center}}

%-- Hyperref ------------------------------------------------------------------
\hypersetup{colorlinks=true,linkcolor=mainblue,urlcolor=mainblue}

%==============================================================================
\begin{document}

\begin{center}
  {\Large\bfseries\color{mainblue} Modos de Cifrado en Bloque}\\[4pt]
  {\normalsize Criptograf\'ia y Seguridad (72.44) --- ITBA \quad|\quad Gu\'ia de estudio para parcial}
\end{center}
\vspace{-4pt}
\hrule
\vspace{8pt}

%==============================================================================
\section{Introducci\'on: por qu\'e se necesitan modos}

Un cifrador de bloque opera sobre bloques de tama\~no fijo $n$ (p.\,ej.\ 128
bits en AES). En la pr\'actica los mensajes suelen ser m\'as grandes (o m\'as
peque\~nos) que un bloque:

\begin{itemize}[nosep]
  \item \textbf{Mensaje m\'as peque\~no que un bloque} --- se aplica
    \emph{padding}. Dos esquemas comunes:
    \begin{itemize}[nosep]
      \item \emph{Simple Pad}: completar con ceros (el receptor debe conocer
        el largo real del mensaje).
      \item \emph{Des Pad} (ISO/IEC 9797-1 M2): agregar un bit 1 y luego
        ceros; puede requerir un bloque adicional completo.
    \end{itemize}
  \item \textbf{Mensaje m\'as grande que un bloque} --- dividir en bloques
    $M_1, M_2, \ldots, M_\ell$, aplicar padding al \'ultimo si necesario y
    transformar cada bloque seg\'un un \textbf{modo de encadenamiento}.
\end{itemize}

\begin{destacado}
  \textbf{Objetivo de un modo}: extender una primitiva de bloque (que solo
  maneja $n$ bits) a mensajes de longitud arbitraria, preservando propiedades
  de seguridad como CPA-Security.
\end{destacado}

Hay varios modos con propiedades distintas; elegir el incorrecto puede romper
la seguridad aunque la primitiva subyacente sea segura.

%==============================================================================
\section{ECB (Electronic Codebook)}

\begin{definicion}{Definici\'on ECB}
  Trata cada bloque de forma \textbf{independiente}:
  \[
    C_i = E_k(M_i), \qquad i = 1, \ldots, \ell.
  \]
  No usa IV ni encadenamiento. El cifrado es \textbf{determinista}.
\end{definicion}

\begin{center}
\begin{tikzpicture}[
  font=\small,
  enc/.style={draw, rounded corners=1pt, minimum width=15mm, minimum height=10mm, fill=blue!6},
  msg/.style={draw, minimum width=15mm, minimum height=7mm},
  arrow/.style={-{Stealth}}
]
  \foreach \i in {1,2,3}{
    \node[msg] (m\i) at (\i*3,0) {$M_{\i}$};
    \node[enc] (e\i) at (\i*3,-1.5) {$E_k$};
    \node[msg] (c\i) at (\i*3,-3) {$C_{\i}$};
    \draw[arrow] (m\i) -- (e\i);
    \draw[arrow] (e\i) -- (c\i);
  }
  % bloque duplicado
  \node[msg, fill=red!15] (m4) at (12,0) {$M_{1}$};
  \node[enc] (e4) at (12,-1.5) {$E_k$};
  \node[msg, fill=red!15] (c4) at (12,-3) {$C_{1}$};
  \draw[arrow] (m4) -- (e4);
  \draw[arrow] (e4) -- (c4);
  \node[font=\footnotesize, text width=2.5cm, align=center] at (12,-4.1)
    {igual a $M_1$\\$\Rightarrow$ cifrado \textbf{igual}};
\end{tikzpicture}
\end{center}
\diagcaption{ECB: cada bloque $M_i$ se cifra independientemente con $E_k$.
  Dos bloques de texto plano iguales producen cifrados iguales (revelando
  patrones).}

\begin{ojo}
  \textbf{ECB NO es CPA-Secure. No utilizar en ning\'un sistema real.}
  Al ser determinista, bloques de texto plano \emph{iguales} $\Rightarrow$
  cifrados \emph{iguales}: se revelan patrones en el mensaje. El ejemplo cl\'asico
  es la imagen del ping\"uino de Linux cifrada con ECB, donde la silueta del
  animal sigue siendo visible. Esto se pregunta en recuperatorios (2025):
  identificar por qu\'e ECB filtra estructura.
\end{ojo}

%==============================================================================
\section{CBC (Cipher Block Chaining)}

\begin{definicion}{Definici\'on CBC}
  Usa la salida del bloque anterior como entrada del siguiente. Requiere un
  \textbf{vector de inicializaci\'on} $C_0 = IV$ elegido aleatoriamente.\\[4pt]
  \textbf{Cifrado:}
  \[
    C_i = E_k\!\left(M_i \oplus C_{i-1}\right), \qquad C_0 = IV.
  \]
  \textbf{Descifrado:} despejando $M_i$ (aplicar $D_k$ y deshacer el XOR):
  \[
    M_i = D_k(C_i) \oplus C_{i-1}.
  \]
\end{definicion}

\subsection*{Diagrama de cifrado CBC}

\begin{center}
\begin{tikzpicture}[
  font=\small,
  enc/.style={draw, rounded corners=1pt, minimum width=14mm, minimum height=9mm, fill=blue!6},
  msg/.style={draw, minimum width=14mm, minimum height=7mm},
  iv/.style={draw, minimum width=14mm, minimum height=7mm, fill=gray!14},
  xor/.style={draw, circle, inner sep=0pt, minimum size=5.5mm},
  arrow/.style={-{Stealth}}
]
  \node[iv] (iv) at (-2.6,-1.6) {$C_0{=}IV$};
  \foreach \i in {1,2,3}{
    \node[msg] (m\i) at (\i*3,0) {$M_{\i}$};
    \node[xor] (x\i) at (\i*3,-1.6) {$\oplus$};
    \node[enc] (e\i) at (\i*3,-3.1) {$E_k$};
    \node[msg] (c\i) at (\i*3,-4.6) {$C_{\i}$};
    \draw[arrow] (m\i) -- (x\i);
    \draw[arrow] (x\i) -- (e\i);
    \draw[arrow] (e\i) -- (c\i);
  }
  \draw[arrow] (iv) -- (x1);
  \draw[arrow] (c1) -| (4.0,-1.6) -- (x2);
  \draw[arrow] (c2) -| (7.0,-1.6) -- (x3);
\end{tikzpicture}
\end{center}
\diagcaption{CBC (cifrado): $C_i=E_k(M_i\oplus C_{i-1})$ con $C_0=IV$. Cada
  bloque de texto plano se XOR-ea con el cifrado anterior antes de pasar por
  $E_k$.}

\subsection*{Diagrama de descifrado CBC}

\begin{center}
\begin{tikzpicture}[
  font=\small,
  dec/.style={draw, rounded corners=1pt, minimum width=14mm, minimum height=9mm, fill=green!8},
  msg/.style={draw, minimum width=14mm, minimum height=7mm},
  iv/.style={draw, minimum width=14mm, minimum height=7mm, fill=gray!14},
  xor/.style={draw, circle, inner sep=0pt, minimum size=5.5mm},
  arrow/.style={-{Stealth}}
]
  \node[iv] (iv) at (-2.6,-3.1) {$C_0{=}IV$};
  \foreach \i in {1,2,3}{
    \node[msg] (c\i) at (\i*3,0) {$C_{\i}$};
    \node[dec] (d\i) at (\i*3,-1.6) {$D_k$};
    \node[xor] (x\i) at (\i*3,-3.1) {$\oplus$};
    \node[msg] (m\i) at (\i*3,-4.6) {$M_{\i}$};
    \draw[arrow] (c\i) -- (d\i);
    \draw[arrow] (d\i) -- (x\i);
    \draw[arrow] (x\i) -- (m\i);
  }
  \draw[arrow] (iv) -- (x1);
  \draw[arrow] (c1) -| (4.0,-3.1) -- (x2);
  \draw[arrow] (c2) -| (7.0,-3.1) -- (x3);
\end{tikzpicture}
\end{center}
\diagcaption{CBC (descifrado): $M_i=D_k(C_i)\oplus C_{i-1}$. El cifrado pasa
  por $D_k$ y luego se deshace el XOR con $C_{i-1}$.}

\begin{destacado}
  CBC con una primitiva pseudoaleatoria (PRP) e IV \textbf{aleatorio} (no
  predecible) es \textbf{CPA-Secure}. Si el IV es predecible, la seguridad se
  rompe (ataque BEAST en TLS 1.0).
\end{destacado}

%==============================================================================
\section{CFB (Cipher Feedback)}

\begin{definicion}{Definici\'on CFB}
  Convierte el cifrador de bloque en un cifrador de flujo, usando \textbf{solo
  $E_k$} (no requiere $D_k$). Opera realimentando el \emph{cifrado} anterior:
  \[
    C_i = E_k(C_{i-1}) \oplus M_i, \qquad C_0 = IV.
  \]
\end{definicion}

\begin{center}
\begin{tikzpicture}[
  font=\small,
  enc/.style={draw, rounded corners=1pt, minimum width=14mm, minimum height=9mm, fill=blue!6},
  blk/.style={draw, minimum width=14mm, minimum height=7mm},
  iv/.style={draw, minimum width=14mm, minimum height=7mm, fill=gray!14},
  xor/.style={draw, circle, inner sep=0pt, minimum size=5.5mm},
  arrow/.style={-{Stealth}}
]
  \node[iv] (iv) at (-2.6,0) {$IV$};
  \foreach \i in {1,2,3}{
    \node[enc] (e\i) at (\i*3,0) {$E_k$};
    \node[xor] (x\i) at (\i*3,-1.6) {$\oplus$};
    \node[blk] (m\i) at (\i*3-1.7,-1.6) {$M_{\i}$};
    \node[blk] (c\i) at (\i*3,-3.0) {$C_{\i}$};
    \draw[arrow] (e\i) -- (x\i);
    \draw[arrow] (m\i) -- (x\i);
    \draw[arrow] (x\i) -- (c\i);
  }
  \draw[arrow] (iv) -- (e1);
  \draw[arrow] (c1) -| (4.0,0.7) -- (4.0,0.0) -- (e2);
  \draw[arrow] (c2) -| (7.0,0.7) -- (7.0,0.0) -- (e3);
\end{tikzpicture}
\end{center}
\diagcaption{CFB: usa solo $E_k$. Se cifra el bloque anterior ($IV$ el
  primero, luego $C_{i-1}$) y la salida se XOR-ea con $M_i$.}

\begin{destacado}
  CFB es \textbf{CPA-Secure} con IV aleatorio. \'Util en dispositivos
  embebidos donde no se puede implementar $D_k$.
\end{destacado}

%==============================================================================
\section{OFB (Output Feedback)}

\begin{definicion}{Definici\'on OFB}
  La \emph{salida} de $E_k$ (keystream) se realimenta a la entrada del
  bloque siguiente, \textbf{sin depender del texto}:
  \[
    S_i = E_k(S_{i-1}), \quad S_0 = IV; \qquad C_i = S_i \oplus M_i.
  \]
\end{definicion}

\begin{center}
\begin{tikzpicture}[
  font=\small,
  enc/.style={draw, rounded corners=1pt, minimum width=14mm, minimum height=9mm, fill=blue!6},
  blk/.style={draw, minimum width=14mm, minimum height=7mm},
  iv/.style={draw, minimum width=14mm, minimum height=7mm, fill=gray!14},
  xor/.style={draw, circle, inner sep=0pt, minimum size=5.5mm},
  arrow/.style={-{Stealth}}
]
  \node[iv] (iv) at (-2.6,0) {$IV$};
  \foreach \i in {1,2,3}{
    \node[enc] (e\i) at (\i*3,0) {$E_k$};
    \node[xor] (x\i) at (\i*3,-1.6) {$\oplus$};
    \node[blk] (m\i) at (\i*3-1.7,-1.6) {$M_{\i}$};
    \node[blk] (c\i) at (\i*3,-3.0) {$C_{\i}$};
    \draw[arrow] (e\i) -- (x\i);
    \draw[arrow] (m\i) -- (x\i);
    \draw[arrow] (x\i) -- (c\i);
  }
  \draw[arrow] (iv) -- (e1);
  \draw[arrow] (e1) -- (4.5,0) -- (e2);
  \draw[arrow] (e2) -- (7.5,0) -- (e3);
\end{tikzpicture}
\end{center}
\diagcaption{OFB: la \emph{salida} de $E_k$ (el keystream $S_i$) se
  realimenta a la entrada del bloque siguiente. A diferencia de CFB, no es
  el cifrado $C_i$ sino el keystream $S_i$ lo que se retroalimenta.}

\begin{destacado}
  \textbf{Clave}: el keystream $S_1, S_2, \ldots$ puede \textbf{precalcularse}
  antes de conocer el texto plano (no depende de $M_i$). OFB es
  \textbf{CPA-Secure} con IV aleatorio.
\end{destacado}

%==============================================================================
\section{CTR (Counter Mode)}

\begin{definicion}{Definici\'on CTR}
  Genera el keystream cifrando un \textbf{contador} (nonce concatenado con un
  \'indice creciente):
  \[
    C_i = E_k(\mathrm{nonce} \| i) \oplus M_i.
  \]
  El descifrado usa exactamente la misma operaci\'on:
  $M_i = E_k(\mathrm{nonce} \| i) \oplus C_i$.
\end{definicion}

\begin{center}
\begin{tikzpicture}[
  font=\small,
  enc/.style={draw, rounded corners=1pt, minimum width=14mm, minimum height=9mm, fill=blue!6},
  ctr/.style={draw, minimum width=18mm, minimum height=8mm, fill=gray!14,
              align=center, font=\scriptsize},
  msg/.style={draw, minimum width=14mm, minimum height=7mm},
  xor/.style={draw, circle, inner sep=0pt, minimum size=5.5mm},
  arrow/.style={-{Stealth}}
]
  \foreach \i [count=\j from 0] in {1,2,3}{
    \node[ctr] (n\i) at (\i*3.3,0) {nonce\,$\|$\,\j};
    \node[enc] (e\i) at (\i*3.3,-1.6) {$E_k$};
    \node[xor] (x\i) at (\i*3.3,-3.1) {$\oplus$};
    \node[msg] (m\i) at (\i*3.3-1.7,-3.1) {$M_{\i}$};
    \node[msg] (c\i) at (\i*3.3,-4.5) {$C_{\i}$};
    \draw[arrow] (n\i) -- (e\i);
    \draw[arrow] (e\i) -- (x\i);
    \draw[arrow] (m\i) -- (x\i);
    \draw[arrow] (x\i) -- (c\i);
  }
\end{tikzpicture}
\end{center}
\diagcaption{CTR: cada bloque es independiente. El keystream se genera
  cifrando \textup{nonce}\,$\|i$; luego se XOR-ea con $M_i$. Totalmente
  paralelizable y con acceso aleatorio. Nunca usa $D_k$.}

\begin{examen}
  \textbf{CTR nunca usa $D_k$ (pregunta 2023).} Como cifrado y descifrado
  aplican la misma operaci\'on ($E_k(\mathrm{nonce}\|i)\oplus\cdot$), CTR es
  v\'alido incluso si la primitiva \textbf{no es invertible}. Lo mismo aplica
  a OFB y CFB. Este es un punto fino que suele preguntarse.
\end{examen}

\begin{destacado}
  CTR es \textbf{CPA-Secure} siempre que no se repita el par
  $(k, \mathrm{nonce})$. Si se reutiliza, XOR-ear dos cifrados revela el XOR
  de los textos planos.
\end{destacado}

%==============================================================================
\section{Tabla resumen de modos}

\begin{center}
\renewcommand{\arraystretch}{1.3}
\begin{tabular}{@{}lcccp{5.2cm}@{}}
\toprule
\textbf{Modo} & \textbf{CPA-Secure} & \textbf{Usa $D_k$} &
\textbf{Paralelizable} & \textbf{Notas clave} \\
\midrule
ECB & \textbf{No} & S\'i & S\'i (enc+dec)
    & Determinista; revela patrones; \textbf{nunca usar} \\
CBC & S\'i (IV aleat.) & S\'i (dec) & Solo descifrado
    & Encadena; IV predecible $\Rightarrow$ inseguro \\
CFB & S\'i (IV aleat.) & No & No (enc); s\'i (dec)
    & Solo $E_k$; flujo a partir de bloque \\
OFB & S\'i (IV aleat.) & No & S\'i (keystream)
    & Keystream precalculable; no depende del texto \\
CTR & S\'i (nonce \'unico) & No & S\'i (enc+dec)
    & Acceso aleatorio; nunca $D_k$; muy eficiente \\
\bottomrule
\end{tabular}
\end{center}

%==============================================================================
\section{Propagaci\'on de errores en CBC (n\'ucleo de examen)}

\begin{examen}
  \textbf{Metodolog\'ia general (aparece casi todos los a\~nos: 2015 PCBC, 2016
  FSM, 2018/2023 CBC, 2025-1, 2025-2):}
  \begin{enumerate}[nosep]
    \item \textbf{Escribir la recurrencia de cifrado} del esquema dado y
      \textbf{despejar $M_i$}: aplicar $D_k$ e invertir los XOR.
    \item \textbf{Analizar un bit de error en $C_j$}: determinar qu\'e bloques
      descifrados quedan corruptos y en qu\'e medida.
    \item \textbf{Analizar borrado/inserci\'on de bloques}: ver qu\'e se
      desincroniza.
  \end{enumerate}
\end{examen}

\subsection{CBC est\'andar: un bit de error en $C_j$}

La ecuaci\'on de descifrado es $M_i = D_k(C_i) \oplus C_{i-1}$. Un bit
erróneo en $C_j$ aparece en dos ecuaciones:

\begin{itemize}[nosep]
  \item $M_j = D_k(\mathbf{C_j}) \oplus C_{j-1}$: el $C_j$ corrupto
    \textbf{pasa por $D_k$} $\Rightarrow$ corrompe \textbf{todo el bloque
    $M_j$} (un bit de entrada de un bloque produce salida completamente
    distinta por el efecto avalancha).
  \item $M_{j+1} = D_k(C_{j+1}) \oplus \mathbf{C_j}$: aqu\'i $C_j$ entra
    por \textbf{XOR directo} $\Rightarrow$ corrompe \textbf{solo el bit
    correspondiente} en $M_{j+1}$.
  \item $M_{j+2}$ y siguientes: $C_j$ ya no aparece $\Rightarrow$
    \textbf{intactos}.
\end{itemize}

\begin{destacado}
  En CBC est\'andar el error es \textbf{autolimitado a 2 bloques}: bloque
  $M_j$ corrupto entero + un bit de $M_{j+1}$.
\end{destacado}

\begin{center}
\begin{tikzpicture}[
  font=\small,
  dec/.style={draw, rounded corners=1pt, minimum width=14mm, minimum height=9mm, fill=green!8},
  msg/.style={draw, minimum width=14mm, minimum height=7mm},
  bad/.style={draw, minimum width=14mm, minimum height=7mm, fill=red!18},
  badall/.style={draw, minimum width=14mm, minimum height=7mm, fill=red!35},
  xor/.style={draw, circle, inner sep=0pt, minimum size=5.5mm},
  arrow/.style={-{Stealth}}
]
  \node[msg]    (cm1) at (0,0)  {$C_{j-1}$};
  \node[bad]    (cj)  at (3,0)  {$C_j^{\;\bigstar}$};
  \node[msg]    (cj1) at (6,0)  {$C_{j+1}$};
  \node[msg]    (cj2) at (9,0)  {$C_{j+2}$};

  \node[dec]    (dj)  at (3,-1.6)  {$D_k$};
  \node[dec]    (dj1) at (6,-1.6)  {$D_k$};
  \node[dec]    (dj2) at (9,-1.6)  {$D_k$};
  \draw[arrow] (cj)  -- (dj);
  \draw[arrow] (cj1) -- (dj1);
  \draw[arrow] (cj2) -- (dj2);

  \node[xor]    (xj)  at (3,-3.0)  {$\oplus$};
  \node[xor]    (xj1) at (6,-3.0)  {$\oplus$};
  \node[xor]    (xj2) at (9,-3.0)  {$\oplus$};
  \draw[arrow] (dj)  -- (xj);
  \draw[arrow] (dj1) -- (xj1);
  \draw[arrow] (dj2) -- (xj2);

  \draw[arrow]       (cm1) |- (1.6,-3.0) -- (xj);
  \draw[arrow, red]  (cj)  -- (3,-0.7) -| (4.6,-3.0) -- (xj1);
  \draw[arrow]       (cj1) |- (7.6,-3.0) -- (xj2);

  \node[badall] (mj)  at (3,-4.4)  {$M_j$};
  \node[bad]    (mj1) at (6,-4.4)  {$M_{j+1}$};
  \node[msg]    (mj2) at (9,-4.4)  {$M_{j+2}$};
  \draw[arrow] (xj)  -- (mj);
  \draw[arrow] (xj1) -- (mj1);
  \draw[arrow] (xj2) -- (mj2);

  \node[font=\scriptsize, align=center] at (3,-5.3)  {todo\\\textbf{corrupto}};
  \node[font=\scriptsize, align=center] at (6,-5.3)  {un\\\textbf{bit}};
  \node[font=\scriptsize, align=center] at (9,-5.3)  {\textbf{intacto}};
\end{tikzpicture}
\end{center}
\diagcaption{Propagaci\'on de error en CBC: un bit erróneo en $C_j$
  ($\bigstar$) pasa por $D_k$ y corrompe \textbf{todo} $M_j$; adem\'as entra
  por XOR en $M_{j+1}$ corrompiendo \textbf{solo ese bit}. Desde $M_{j+2}$
  todo queda intacto.}

\subsection{PCBC (Propagating CBC) --- parcial 2015}

En PCBC el encadenamiento incluye tambi\'en el texto plano anterior:
\[
  C_i = E_k(M_i \oplus C_{i-1} \oplus M_{i-1}).
\]
Despejando:
\[
  M_i = D_k(C_i) \oplus C_{i-1} \oplus M_{i-1}.
\]
Un error en $C_j$ afecta $M_j$ (pasa por $D_k$); el $M_j$ corrupto entra
como $M_{i-1}$ en la ecuaci\'on de $M_{j+1}$, que a su vez contamina
$M_{j+2}$, y as\'i sucesivamente. El error \textbf{se propaga a todos los
bloques siguientes} (nunca se autolimita).

\begin{ojo}
  \textbf{PCBC} es la variante \emph{no autolimitada}: un \'unico bit de
  error en cualquier $C_j$ destruye el descifrado desde $j$ en adelante.
  Contrastar con CBC est\'andar (solo 2 bloques afectados).
\end{ojo}

\subsection{Borrado e inserci\'on de bloques en CBC}

\begin{itemize}[nosep]
  \item \textbf{Borrar $C_0 = IV$}: como $M_1 = D_k(C_1)\oplus C_0$, se
    pierde la referencia para $M_1$ (desincronizaci\'on); los dem\'as bloques
    se descifran con un corrimiento de \'indice.
  \item \textbf{Borrar el \'ultimo bloque $C_\ell$}: se pierde \textbf{solo
    $M_\ell$}; todos los bloques anteriores quedan intactos.
  \item \textbf{Insertar un bloque falso}: el bloque falso y el siguiente
    quedan corruptos; desde ahí el descifrado se recupera.
\end{itemize}

\begin{ojo}
  El error m\'as com\'un en estos ejercicios es no distinguir el camino del
  bloque corrupto: \textbf{pasar por $D_k$ corrompe el bloque entero}
  (efecto avalancha); \textbf{entrar por XOR corrompe solo ese bit}.
  Siempre contabilizar en cuntos bloques aparece $C_j$ en las ecuaciones
  de descifrado.
\end{ojo}

%==============================================================================
\section{Memory hooks: mnem\'onicos para el parcial}

\begin{center}
\renewcommand{\arraystretch}{1.4}
\begin{tabular}{@{}ll@{}}
\toprule
\textbf{Modo} & \textbf{Regla mnem\'onica} \\
\midrule
ECB & \textbf{E}stupido --- \textbf{C}opia tu \textbf{B}loque: igual entrada = igual salida. \\
CBC & \textbf{C}adena de \textbf{B}loques \textbf{C}ompartidos: XOR con el anterior. \\
CFB & \textbf{C}ifrado \textbf{F}eedback: el \textbf{C}ifrado alimenta al siguiente. \\
OFB & \textbf{O}utput: el keystream se \textbf{O}btiene antes del texto (precalculable). \\
CTR & \textbf{C}on\textbf{T}ado\textbf{R}: nonce + contador, todo paralelo, nunca $D_k$. \\
\bottomrule
\end{tabular}
\end{center}

\begin{itemize}[nosep]
  \item \textbf{``Solo $E_k$''}: CFB, OFB, CTR. (CBC y ECB necesitan $D_k$
    para descifrar.)
  \item \textbf{Paralelizable en cifrado}: ECB, CTR. (CBC, CFB, OFB no,
    por dependencia secuencial.)
  \item \textbf{Keystream independiente del texto}: OFB y CTR. (CFB y CBC
    dependen del texto cifrado anterior.)
  \item \textbf{Error autolimitado en CBC}: 2 bloques (el $j$ entero + 1
    bit del $j+1$). PCBC: ilimitado.
\end{itemize}

%==============================================================================
\section{Ejercicios tipo examen con respuesta}

\begin{examen}
\textbf{Ejercicio 1 --- V/F/Justificar (estilo 2015 PCBC).}\\
``En PCBC, si se produce un error de un bit en $C_3$, entonces $M_4$ queda
correctamente descifrado.''

\medskip
\textbf{Respuesta: FALSO.} La ecuaci\'on de descifrado PCBC es
$M_i = D_k(C_i)\oplus C_{i-1}\oplus M_{i-1}$. Un error en $C_3$ corrompe
$M_3$ (pasa por $D_k$, efecto avalancha). El $M_3$ corrupto entra como
$M_{i-1}$ en la ecuaci\'on de $M_4$, corrompiendo $M_4$ tambi\'en. El error
se propaga indefinidamente: $M_3, M_4, M_5, \ldots$ quedan todos corruptos.
\end{examen}

\begin{examen}
\textbf{Ejercicio 2 --- V/F/Justificar (estilo 2023 CBC con error).}\\
``En CBC est\'andar, si hay un error de un bit en $C_5$, entonces $M_7$ se
descifra correctamente.''

\medskip
\textbf{Respuesta: VERDADERO.} La ecuaci\'on es $M_i = D_k(C_i)\oplus
C_{i-1}$. El bloque corrupto $C_5$ aparece como:
\begin{itemize}[nosep]
  \item Entrada a $D_k$ en $M_5$: todo $M_5$ corrupto.
  \item XOR directo en $M_6$: solo el bit correspondiente de $M_6$ corrupto.
  \item $C_5$ \textbf{no aparece} en la ecuaci\'on de $M_7 = D_k(C_7)\oplus
    C_6$: $M_7$ intacto.
\end{itemize}
El error se autolimita a 2 bloques: $M_5$ (entero) y $M_6$ (1 bit).
\end{examen}

\begin{examen}
\textbf{Ejercicio 3 --- V/F/Justificar (estilo 2025 CTR).}\\
``Para descifrar con CTR se necesita implementar la funci\'on de descifrado
$D_k$ del cifrador de bloque subyacente.''

\medskip
\textbf{Respuesta: FALSO.} En CTR el descifrado es
$M_i = E_k(\mathrm{nonce}\|i)\oplus C_i$, id\'entico al cifrado. \textbf{Solo
se usa $E_k$}. Esto significa que CTR (al igual que OFB y CFB) es v\'alido
incluso si la primitiva de bloque subyacente no es invertible. No se requiere
ni implementa $D_k$ en ninguno de los modos de flujo (CFB, OFB, CTR).
\end{examen}

\begin{examen}
\textbf{Ejercicio 4 --- Pregunta abierta (propagaci\'on en CBC).}\\
En un mensaje de 5 bloques cifrado con CBC, el tercer bloque de cifrado
($C_3$) se corrompe durante la transmisi\'on. \textbf{Indicar cu\'ales bloques
de texto plano se descifran incorrectamente y por qu\'e.}

\medskip
\textbf{Respuesta.} Usando $M_i = D_k(C_i)\oplus C_{i-1}$:
\begin{itemize}[nosep]
  \item $M_3 = D_k(C_3)\oplus C_2$: $C_3$ corrupto pasa por $D_k$
    $\Rightarrow$ \textbf{$M_3$ entero corrupto}.
  \item $M_4 = D_k(C_4)\oplus C_3$: $C_3$ corrupto entra por XOR
    $\Rightarrow$ \textbf{un bit de $M_4$ corrupto}.
  \item $M_1, M_2, M_5$: \textbf{intactos}.
\end{itemize}
Conclusi\'on: bloques afectados son \textbf{$M_3$ (completo) y $M_4$ (1 bit)}.
\end{examen}

\end{document}
